\begin{solution}
In the inertial frame the light angular speeds are \(\pm1/R\), whereas the
receiver moves with \(\Omega\).  Thus
\[
t_\pm=\frac{2\pi R}{1\mp\Omega R},\qquad
\Delta t=\frac{4\pi\Omega R^2}{1-\Omega^2R^2}.
\]
The ring clock runs slower by \(\sqrt{1-\Omega^2R^2}\), so
\[
\Delta\tau=\frac{4\pi\Omega R^2}{\sqrt{1-\Omega^2R^2}}
\simeq4\Omega A.
\]
Rotation has nonzero vorticity, so no single inertial frame is comoving with
the entire loop.
\end{solution}

\begin{solution}
The required derivative is
\[
\frac{DS^\mu}{D\tau}
=(u^\mu a_\nu-a^\mu u_\nu)S^\nu.
\]
Contracting with \(u_\mu\), using \(u^2=-1\) and \(a\cdot u=0\), shows that
\(\dd(S\cdot u)/\dd\tau=0\).  The derivative contains only the boost needed
to follow the changing \(u^\mu\); its spatial projection in the momentary rest
frame vanishes, so there is no imposed rotation.
\end{solution}

\begin{solution}
Geodesic deviation gives
\[
\frac{D^2\xi^i}{D\tau^2}=-R^i{}_{0j0}\xi^j.
\]
Near a spherical mass \(R_{0\ii0j}\sim M/r^3\).  Across size \(L\), the
differential acceleration is \(\Delta a\sim ML/r^3\), while the mean
acceleration is \(M/r^2\).  Thus \(\Delta a/a\sim L/r\); a uniform elevator
requires both \(L/r\ll1\) and curvature effects \(ML^2/r^3\ll1\).
\end{solution}

\begin{solution}
A static observer has four-velocity
\(u=(1+a\xi)^{-1}\partial_\eta\).  The conserved photon quantity
\(-k\cdot\partial_\eta\) gives
\[
\frac{\nu_2}{\nu_1}=\frac{1+a\xi_1}{1+a\xi_2}.
\]
For \(a|\xi_2-\xi_1|\ll1\),
\(\Delta\nu/\nu=-a(\xi_2-\xi_1)\), the weak-field result
\(\Delta\nu/\nu=-\Delta\Phi\).  The lapse vanishes at
\(\xi=-a^{-1}\), the Rindler horizon.
\end{solution}

\begin{solution}
The accelerations are \(a_A=g(1+\delta_A)\) and
\(a_B=g(1+\delta_B)\).  Hence, to first order,
\[
\eta_{AB}
=\frac{2|\delta_A-\delta_B|}{2+\delta_A+\delta_B}
\simeq|\delta_A-\delta_B|.
\]
The experiment tests universality of free fall, or the weak equivalence
principle, by looking for composition dependence.
\end{solution}

\begin{solution}
Along the unperturbed path \(r^2=b^2+x^2\), the transverse acceleration is
\(a_\perp=Mb/r^3\).  Therefore
\[
\Delta\theta_{\rm EP}
=\int_{-\infty}^{\infty}\frac{Mb\,\dd x}{(b^2+x^2)^{3/2}}
=\frac{2M}{b}.
\]
General relativity gives \(4M/b\).  The elevator argument captures the
time-time metric perturbation but misses the equal contribution from spatial
curvature.
\end{solution}

\begin{solution}
Rindler worldlines at fixed \(\xi\) have constant mutual proper distance on
\(\eta=\) constant slices.  Their lapse is \(N=1+a_0\xi\), and their
four-acceleration magnitude is
\[
a(\xi)=\partial_\xi\ln N=\frac{a_0}{1+a_0\xi}.
\]
Equal proper acceleration would make initially separated hyperbolae change
their proper spacing; the rear must accelerate more strongly than the front.
\end{solution}

\begin{solution}
For \(\dd s^2\simeq-(1+2\Phi)\dd t^2+\dd\bm x^2\),
\(\dd\tau\simeq(1+\Phi)\dd t\).  Since
\(\Delta\Phi\simeq gh\),
\[
\Delta\tau\simeq gh\,T
\]
in units \(c=1\), or \(ghT/c^2\) in SI units.  For \(h=1\,\mathrm m\) and one
day, this is about \(9\times10^{-12}\,\mathrm s\), several picoseconds.
\end{solution}

\begin{solution}
With \(\phi'=\phi-\Omega t\),
\[
\dd s^2=-(1-\Omega^2r^2)\dd t^2
+2\Omega r^2\dd t\,\dd\phi'
+\dd r^2+r^2\dd\phi'^2+\dd z^2.
\]
A stationary rotating clock has
\(\dd\tau=\sqrt{1-\Omega^2r^2}\,\dd t\).  Expanding the geodesic equation at
low speed yields
\[
\ddot{\bm r}_{\rm rot}
=-2\bm\Omega\times\bm v_{\rm rot}
-\bm\Omega\times(\bm\Omega\times\bm r),
\]
the Coriolis and centrifugal terms.
\end{solution}

\begin{solution}
For masses initially at rest,
\[
\ddot\xi^i=-E^i{}_j\xi^j.
\]
Applying three linearly independent separation vectors determines the three
columns of \(E^i{}_j\).  Symmetry provides a consistency check.  In vacuum
\(E^i{}_i=R_{00}=0\), leaving five independent components; suitably chosen
accelerations along three independent directions still suffice to reconstruct
the full trace-free symmetric matrix.
\end{solution}
